% Generated draft: requires mathematical and source review.
\documentclass[11pt]{article}
\usepackage[margin=25mm]{geometry}
\usepackage{fontspec}
\setmainfont{Latin Modern Roman}
\usepackage{amsmath,amssymb,mathtools,unicode-math}
\setmathfont{Latin Modern Math}
\usepackage{enumitem,needspace,xurl,hyperref}
\setlength{\parindent}{0pt}
\setlength{\parskip}{5pt}
\setlength{\emergencystretch}{3em}
\newcommand{\N}{\mathbb{N}}
\newcommand{\Z}{\mathbb{Z}}
\newcommand{\Q}{\mathbb{Q}}
\newcommand{\R}{\mathbb{R}}
\newcommand{\C}{\mathbb{C}}
\newcommand{\F}{\mathbb{F}}
\newcommand{\A}{\mathbb{A}}
\newcommand{\PP}{\mathbb{P}}
\newcommand{\E}{\mathbb{E}}
\newcommand{\eps}{\varepsilon}
\newcommand{\dd}{\,\mathrm{d}}
\begin{document}
\section*{UnsolvedMath ID: 20002035}
\textbf{Source problem code:} catalogue entry
\subsection{Definitions and mathematical statement}
The isolated source sentence does not define $t(g)$. Following the supplied aid, interpret it as the complex dimension of the global twistor-spinor space for a \emph{chosen} spin structure on a connected conformal spin manifold $(M^n,[g])$, $n\ge3$, of signature $(p,q)$. With Clifford convention $X\cdot X=-g(X,X)$, a twistor spinor satisfies
\[
 \nabla_X\psi+\frac1nX\cdot\mathcal D_g\psi=0
 \quad\text{for every }X\in TM,
 \qquad t(g)=\dim_\C\ker\mathcal P_g,
\]
where $\mathcal P_g$ is this twistor operator and $\mathcal D_g$ the metric Dirac operator. These definitions are an explicit interpretation of the missing antecedent, not wording present in Problem 15.

\textbf{Problem 15.} Relate $t(g)$ to the holonomy of conformal Cartan connections.

\textbf{Supplied spin-tractor formulation.} Let $\mathcal S_{\mathrm{tr}}$ be the complex spin-tractor bundle for the chosen conformal spin lift of the normal conformal Cartan connection. The aid invokes the correspondence between twistor spinors and parallel spin tractors. At a basepoint, let $\Sigma$ be its spin representation space and $H\subset\operatorname{Spin}(p+1,q+1)$ the actual full spin-tractor holonomy. Then the asserted exact formula is
\[
 t(g)=\dim_\C\Sigma^H.
\]
The ordinary vector holonomy in an orthogonal group does not specify the spin lift or its global signs and is insufficient to determine this dimension alone. The formula uses full holonomy, not only its Lie algebra or restricted subgroup.

Let $H^0$ denote restricted spin-tractor holonomy and $W=\Sigma^{H^0}$. On the universal cover, the twistor-spinor space is identified with $W$. The remaining monodromy gives a representation $\rho:\pi_1(M)\to\operatorname{GL}(W)$, and
\[
 t(g)=\dim_\C W^{\rho(\pi_1(M))}.
\]
If the spin structure is twisted by $\chi\in H^1(M;\Z_2)=\operatorname{Hom}(\pi_1(M),\{\pm1\})$, its dimension becomes
\[
 t_\chi=\dim_\C\{w\in W:\rho(\gamma)w=\chi(\gamma)w
                                  \text{ for all }\gamma\in\pi_1(M)\}.
\]
Here multiplicity means $\dim\operatorname{Hom}_{\pi_1(M)}(\C_\chi,W)$, without a semisimplicity assumption for an infinite fundamental group. If $\Gamma=\pi_1(M)$ is finite, this specializes to
\[
 t_\chi=\frac1{|\Gamma|}\sum_{\gamma\in\Gamma}
                      \chi(\gamma)^{-1}\operatorname{tr}\rho(\gamma).
\]

\textbf{Supplied flat diagnostic.} On a flat Lorentzian $n$-torus, write $r=2^{\lfloor n/2\rfloor}$ for the rank of the full complex base spinor bundle. The aid states that the universal cover has $2r$ twistor spinors, the trivial torus spin structure has $r$ global ones, and every nontrivial spin structure has zero. Restricted flat holonomy alone would miss the global translation monodromy and spin signs. These are the aid's claimed examples supporting the full-holonomy distinction, not new computations or a classification of all conformal holonomy groups.
\subsection{Short English statement}
Interpreting the missing $t(g)$ as the twistor-spinor dimension, the supplied aid identifies it with the fixed-spinor dimension of full spin-tractor holonomy. Restricted holonomy describes the universal cover; fundamental-group monodromy and the chosen spin structure determine which solutions descend globally. Vector conformal holonomy alone omits this information.
\subsection{Sources}
\begin{itemize}
\item Local UnsolvedMath registry, permanent ID 20002035, source code catalogue entry; canonical \texttt{aim-geometry-notes.json}, notes[372].
\item AIM \emph{Conformal structure in geometry, analysis, and physics}, Problem 15: \url{https://aimath.org/WWN/confstruct/confstruct.pdf}. The supplied verification reports the official PDF version dated 15 October 2003.
\item No separate bibliographic citation defining $t(g)$ or proving the twistor-spinor/parallel-spin-tractor correspondence and monodromy formulas is supplied. Their interpretation and claimed results come from the registry research aid.
\item Registry research attempt 1, dated 2026-08-14: full solution, problem status at run solved, rigor ``proved'', novelty confidence low. The explicitly machine-generated, non-peer-reviewed aid requires independent verification before citation.
\end{itemize}
\end{document}
